\centerline{\bf \Huge Двойные двойные отношения}

\begin{flushright}
        \scriptsize{последний листочек этой смены(}
\end{flushright}

\problem{1}
Прямые $a, b, c, d$ пересекаются в одной точке. Известно, что $c \perp d$ и  $(a, b ; c, d)=-1$. Тогда $c$ и $d$ - биссектрисы углов, образованных прямыми $a$ и $b$.

\problem{2}
Диагонали выпуклого четырёхугольника $A B C D$ пересекаются в точке $R$, продолжения сторон $A B$ и $C D$ - в точке $P$, а продолжения сторон $B C$ и $A D$ в точке $Q$. Прямая, проходящая через точку $R$, пересекает прямые $A B, B C, C D$, $D A$ в точках $M, X, N, Y$. Докажите, что $R$ делит пополам отрезок $M N$ тогда и только тогда, когда $R$ делит пополам отрезок $X Y$.

\problem{3}
В обозначениях предыдущей задачи пусть $T$ --- проекция точки $R$ на прямую $P Q$. Докажите, что противоположные стороны четырёхугольника видны из точки $T$ под равными углами.

\problem{4}
\textbf{Окружность Аполлония.} Даны точки $A$ и $B$ и действительное $k>1$. Докажите, что ГМТ $X$, для которых $A X / B X=k$ является окружность.

\problem{5}
Внутри угла с вершиной $O$ отмечена точка $P$. Рассматриваются всевозможные пары точек $X$ и $Y$ на сторонах угла такие, что $\angle O P X=\angle O P Y$. Докажите, что все прямые $X Y$ проходят через одну точку.

\problem{6}
В треугольнике $A B C$ проведена биссектриса $A L$ и биссектриса внешнего угла $A N$. Точка $M$ - середина стороны $A C, P$ - точка пересечения прямых $A B$ и $M L$. Докажите, что $P A=P N$.

\problem{7}
Вписанная окружность касается сторон $A B, A C, B C$ треугольника $A B C$ в точках $C_1, B_1, A_1$ соответственно. Точка $H$ - проекция $A_1$ на прямую $B_1 C_1$. Докажите, что $H A_1$ - биссектриса угла $B H C$.

\problem{8}
Пусть $H_a$ и $L_a$ - основание высоты и биссектрисы треугольника $A B C$, проведённых из вершины $A$; вписанная и вневписанная окружности касаются стороны $B C$ в точках $K_a$ и $T_a$ соответственно.

(a) Докажите, что $\left(H_a, L_a ; K_a, T_a\right)=-1$;


\newpage

\hspace{3mm}

(б) Пусть $I$ и $I_a$ - центр вписанной и вневписаной окружностей треугольника $A B C$. Докажите, что прямые $H_a I$ и $H_a I_a$ симметричны относительно прямой $B C$.

(в) Определим аналогично точки $H_b, L_b, K_b, T_b$. Докажите, что прямые $H_a H_b$, $L_a L_b$, $T_a T_b$, $K_a K_b$ пересекаются в одной точке.

\problem{9}
Дан выпуклый четырехугольник $A B C D$. Точки $Q, A, B, P$ лежат на одной прямой именно в таком порядке. Прямая $A C$ касается окружности $(A D Q)$, а прямая $B D$ касается окружности ( $B C P$ ). Точки $M$ и $N$ - середины сторон $B C$ и $A D$ соответственно. Докажите, что следующие касательная к $(A N Q)$ в точке $A$, касательная к ( $B M P$ ) в точке $B$ пересекаются на прямой $C D$.

\problem{10}
На стороне $B C$ треугольника $A B C$ выбрана точка $D$. Окружности, вписанные в треугольники $A B D$ и $A C D$ касаются отрезка $B C$ в точках $E$ и $F$ соответственно. Линия центров окружностей пересекает отрезок $A D$ в точке $P$. Прямые $B P$ и $C P$ пересекают биссектрисы углов $C$ и $B$ в точках $Y$ и $X$ соответственно. Докажите, что прямые $E X$ и $F Y$ пересекаются на вписанной окружности треугольника $A B C$.